\section{Limitations and open problems}
\label{sec:limits}

\paragraph{No token is reused inside a model.} Every result here is about which patches a rule sends, judged in pixels. Whether a video LLM fed only the sent patches answers as it does on full frames is measured in \citet{nair2026certifiedskip} on fixed cameras with a rebuilt-image protocol, and not at all on moving cameras. Doing it properly needs a model whose encoder accepts a sparse set of patches with their positions, plus the warp as metadata on frozen tokens. We have not built that.

\paragraph{No long-video benchmark.} There is no result on Video-MME \citep{fu2024videomme}, LongVideoBench \citep{wu2024longvideobench} or MLVU \citep{zhou2024mlvu}. Those benchmarks measure answer accuracy at a frame budget, and most of their footage is edited video at 1--2~fps, where a cut or a large change every few frames is the norm. The audit of Section~\ref{sec:audit} suggests what to expect there: at 2~fps the consecutive rules recover some slow change, and the advantage of anchoring shrinks to the events that a low rate makes look brief.

\paragraph{No certificate through a network.} The guarantee is on pixels. Section~\ref{sec:shift} explains why it does not pass through a standard encoder: dot-product attention is not Lipschitz on an unbounded domain, so the product of layer constants is undefined or vacuous. A certificate on tokens would need either an encoder built with bounded Lipschitz constants or a bound that uses the structure of real activations, and we claim neither.

\paragraph{The interpolation term is large on edges.} Theorem~\ref{thm:warp} is exact relative to the constructed reference $W$, but $W$ is a resampling of what was kept, and $\iota$ measures how far apart the two can be. On DAVIS its median over dropped patches is tens of grey levels, and on sharp edges it reaches the edge's contrast. Where $\iota$ is large, the statement relative to the kept pixels is weak, and Proposition~\ref{prop:iota} is honest about that rather than a fix for it. A reference built only from whole-pixel moves would make $\iota$ zero at the cost of keeping more under sub-pixel motion. The synthetic audit (Section~\ref{sec:warpaudit}) shows no loss of $\Delta$-features in number, which is weaker than a proof.

\paragraph{The competitive bound is for a fixed warp.} Theorem~\ref{thm:opt} holds for any seminorm on a fixed camera, and for the warp rule with the warp held fixed. When the warp changes between frames, the triangle inequality picks up an interpolation term and we state no bound. The measured ratios of Section~\ref{sec:compexp} are on fixed cameras only.

\paragraph{Speed.} The warp rule is a Python prototype and took about 1.0--1.4~s per 480p frame for registration, refinement and the canvas (nine clips in parallel on one laptop), against about 0.15~s for the fixed-camera rule in Python under the same load and hundreds of frames per second for its C++ core. Homography registration is cheap in principle and runs in real time in every video codec; the gap is engineering, not mathematics, but it has not been closed.

\paragraph{The audit tests decision rules, not systems.} Each reducer is reduced to its temporal comparison in pixel space (Appendix~\ref{app:reducers}). A learned feature can be more or less sensitive than the pixel statistic that replaces it, so the recall numbers are not those of the published systems. What carries over is the structural point: no feature, learned or not, sees a change against a frame it is not compared with. DyCoke's window of four frames was implemented but not run in this audit, PruneVid's clustering into segments and LongVU's query-guided and spatial steps are not re-implemented, and the attention-based reducers (FastV, VisionZip, SparseVLM) need the model and are out of scope.

\paragraph{Footage.} Fixed-camera results use three clips (hall monitor, a CAVIAR lobby, one VIRAT scene) for the audit, the competitive analysis and the shift cost; moving-camera results use DAVIS~2017 at 40 frames per clip and three excerpts of one film. The planted events are synthetic objects composited into real footage. A deployment would want a larger and more varied set of cameras, and real events with ground truth.
